% ============================================================
% Chapter 6: Experimental Design
%   Section: Model Benchmarks
%     Subsection: Writing Models
% Included from chapters/06-experimental-design/06-model-benchmarks/model-benchmarks.tex
% Source: docs/decisions/evaluation.md, "Writing-model benchmark";
% backend/src/evaluation/writing.py; data/evaluation/writing/.
% Moved to P2 ("only if time is left") on 14 October: it answers none
% of the four research questions.
% ============================================================

\subsection{Writing Models}
\label{subsec:benchmark-writing}

The text corrector (Section~\ref{sec:design-interface}) judges a text on
seven metrics; two of them, readability and coverage of verification,
are computed deterministically, and five --- grammar, factual
consistency, search-engine optimisation, hallucination and style --- come
from one call to the language model. The benchmark measures a model as
the \emph{editor} it is in that role, by running the corrector's own
call, prompt included, and refusing to compare runs made with different
prompts. It answers none of the four research questions and was moved
to ``only if time is left'' on 14~October; it is described here because
it is built and tested.

\paragraph{Planted defects, not opinions.} A human score for how good a
text's search optimisation is would be an opinion per text; a planted
defect is a fact. Six short news texts were written for the set, three
in English and three in Spanish, each in a clean version and five
versions with exactly one defect aimed at one metric: agreement,
spelling and tense errors (grammar); a figure contradicting an earlier
one (factual consistency); an invented authority and sweeping claims
(hallucination); a vague headline with the subject gone from the lead
(search optimisation); and a casual, repetitive, exclamatory register
(style). Each defect version against its clean text is a pair whose
right answer is known --- the targeted score should go down --- which
makes 30 pairs, 15 per language.

\paragraph{The rubric}, in the order a model is judged:
\begin{enumerate}
    \item \textbf{Format compliance} as a gate: the share of texts for
    which all five metrics came back usable. A model below 95\% is out,
    whatever it scores when it complies; a missing or non-numeric
    score is never averaged as zero.
    \item \textbf{Detection}: per pair, 1 when the defect version scored
    lower on the targeted metric, 0.5 on a tie, 0 when higher, so a
    model that gives every text the same score sits at chance (0.5);
    with a bootstrap interval over pairs, per metric and per language.
    \item \textbf{Defect named}: whether the targeted metric's summary
    names the defect.
    \item \textbf{Off-target drift}: the mean change on the four metrics
    a defect was not aimed at; lower is better.
    \item \textbf{Consistency}: the spread of a text's score across
    repeats.
    \item \textbf{Calibration}, read rather than ranked: the mean score
    of the clean texts.
    \item \textbf{Cost and latency}, to choose between models the
    rubric cannot separate.
\end{enumerate}

The local baseline is compared with the other local candidates, and
with one hosted provider if time allows, each paired against the
baseline over the pairs both scored. With 30 pairs an interval of about
$\pm15$ points is expected, so only large differences will show. Two
defects of the corrector were found while building the benchmark and
fixed: a null list of issues lost all five metrics, and a missing score
counted as a judgement of zero.

\paragraph{The draft reviewer, on seeded drafts.} The benchmark judges
free text; the reviewer of the writing flow
(Section~\ref{sec:dev-writing}) judges a draft against the sources it
was written from, and proposes a correction that is offered only if it
passes the checks. It is tried apart, on the two drafts seeded with
defects on 15~October, one English and one Catalan (the Catalan one as
a test keeps it, the English one rebuilt from the record of that day),
with the same two sources and fact-checks: in English a spelling error (``it's'' for
``its''), a figure the cited source contradicts (3 million euros where
it says 1.2), a claim the fact-checks judged \texttt{FALSE} stated
uncited, an eight-character title and no subtitle; in Catalan a spelling
error (``esta'' for ``està'') and the same \texttt{FALSE} claim. A
defect counts as found when an issue the \emph{model} raised names it,
read from its answer before the deterministic checks bind the scores;
each draft is reviewed twice per model. Counted too: the seconds and
tokens per call, whether every metric parsed, and whether the correction
was offered or why not. Issues the model invented are read by a person.
The script is \texttt{scripts/try\_reviewer\_model.py}.
